% =============================================================================
%  Hierarchical Bayesian growth model, and comparison with the hidden Markov
%  model from Weeks 5-7.
%
%  Figures from outputs/model/figures and outputs/hbayes/figures, regenerated by
%      python scripts/03_progression_model.py
%      python scripts/04_hierarchical_bayes.py
% =============================================================================
\documentclass[11pt,a4paper]{article}

\usepackage[T1]{fontenc}
\usepackage[utf8]{inputenc}
\usepackage[margin=2.4cm]{geometry}
\usepackage{booktabs}
\usepackage{array}
\usepackage{xcolor}
\usepackage{graphicx}
\usepackage{enumitem}
\usepackage{caption}
\usepackage{amsmath}
\usepackage{amssymb}
\usepackage{float}
\usepackage[hidelinks]{hyperref}

\graphicspath{{../outputs/hbayes/figures/}{../outputs/model/figures/}}

\newcommand{\code}[1]{\texttt{\small #1}}
\newcommand{\arr}{$\rightarrow$}

\setlist[itemize]{itemsep=2pt,topsep=4pt}
\setlist[enumerate]{itemsep=2pt,topsep=4pt}
\captionsetup{font=small,labelfont=bf}

\newenvironment{keypoint}
  {\begin{center}\begin{minipage}{0.94\textwidth}%
   \vspace{4pt}\hrule\vspace{7pt}\small}
  {\vspace{7pt}\hrule\vspace{4pt}\end{minipage}\end{center}}

\title{\textbf{The Hierarchical Bayesian Growth Model}\\[0.4em]
       \large Implementation, and comparison with the hidden Markov model}
\author{Parv\\[0.3em]
        \small Supervisor: Prof.\ K.~R.\ Sahasranand\\
        \small Department of Data Science, IIT Palakkad}
\date{August 2026}

\begin{document}
\maketitle

\begin{abstract}
\noindent
The hierarchical Bayesian growth model specified in the proposal (pages 13--15)
has been implemented as written and compared against the hidden Markov model
from Weeks 5--7. Neither replaces the other; they answer different questions,
and the comparison is informative in both directions.

On predicting the next recorded subtype, the hidden Markov model is ahead in
\textbf{6 of 6} independent splits, by $+0.058 \pm 0.019$ log-likelihood per
transition. The reason is specific and measurable: the growth model's residual
standard deviation is $\sigma = 1.138$, which is \textbf{4.6 times the width of
the entire Indeterminate band}, so band membership is dominated by noise.

The growth model nevertheless supplies something no Markov model can: an
interpretable progression rate on the clinical scale. The TD/PIGD ratio
declines by $\mathbf{0.093}$ per year (95\% credible interval $0.068$ to
$0.116$), with a between-patient standard deviation of $0.103$ per year --- the
individualised trajectory the proposal is ultimately after.

Two problems with the model as specified were found by implementing it. The
observation and process variances are provably not separable, and a Normal
model on the raw ratio places 12.2\% of its five-year predictive mass on
negative ratios, which are impossible.
\end{abstract}

% =============================================================================
\section{The model, as specified}
% =============================================================================

From the proposal, pages 13--15:
\begin{align}
  R^{obs}_{it} &= R_{it} + \epsilon^{obs}_{it},
    &\epsilon^{obs}_{it} &\sim \mathcal{N}(0, \sigma^2_{obs}) \\
  R_{it} &= (\beta_0 + b_{0i}) + (\beta_1 + b_{1i})\, t + \epsilon^{proc}_{it},
    &\epsilon^{proc}_{it} &\sim \mathcal{N}(0, \sigma^2_{proc}) \\
  b_{0i} &\sim \mathcal{N}(0, \tau_0^2), &b_{1i} &\sim \mathcal{N}(0, \tau_1^2)
\end{align}
with the clinical translation
\[
  P(\mathrm{TD}) = P(R > 1.15), \quad
  P(\mathrm{Ind}) = P(0.90 \leq R \leq 1.15), \quad
  P(\mathrm{PIGD}) = P(R < 0.90).
\]

Inference is by Gibbs sampling with conjugate updates throughout
(\code{src/ppmi\_subtype/hbayes.py}); no external probabilistic-programming
dependency. A full run over 573 patients takes about one second. The sampler is
validated by simulating trajectories from chosen parameters and checking the
posterior covers them (\code{tests/test\_hbayes.py}).

\subsection{A note on the two variance terms}

Substituting (2) into (1):
\[
  R^{obs}_{it} = (\beta_0 + b_{0i}) + (\beta_1 + b_{1i})\, t
    + \underbrace{(\epsilon^{proc}_{it} + \epsilon^{obs}_{it})}_{\text{variance } \sigma^2_{proc} + \sigma^2_{obs}}
\]

\begin{keypoint}
Both error terms are independent and additive on the \emph{same} observation,
so the likelihood depends on them only through their sum. Holding
$\sigma^2_{proc} + \sigma^2_{obs}$ fixed and moving mass between the two leaves
the log-likelihood numerically unchanged --- verified directly. They are
\textbf{not separately identifiable} in this parameterisation.

The sampler therefore estimates the identified quantity
$\sigma^2 = \sigma^2_{proc} + \sigma^2_{obs}$ rather than reporting a split the
data cannot support. This matters for the headline output: $P(TD) = P(R^{true}
> 1.15)$ refers to the \emph{latent} ratio, whose spread excludes observation
noise. Without separating the two, only the predictive distribution of
$R^{obs}$ is identified.

Separating them needs either the state-space form the proposal also mentions
(where process noise accumulates over time and so leaves a distinguishable
signature) or repeated measurements at a single visit.
\end{keypoint}

% =============================================================================
\section{Data}
% =============================================================================

\begin{table}[H]
\centering
\begin{tabular}{@{}lr@{}}
\toprule
Labelled visits                          & 4{,}441 \\
With a finite TD/PIGD ratio              & 3{,}716 \quad (83.7\%) \\
Without (PIGD subscore is zero)          & \phantom{0,}725 \\
\midrule
Training observations / patients         & 2{,}426 / 573 \\
\bottomrule
\end{tabular}
\caption{The growth model is defined on the ratio, so the 725 visits with a
zero PIGD subscore cannot enter it. They are 665 TD and 60 Indeterminate, so
the loss is not random --- it removes the most tremor-dominant visits.}
\end{table}

\noindent
The Markov models of Weeks 5--7 have no such restriction, which is why the
comparison in Section~4 is confined to transitions both models can score.

% =============================================================================
\section{Fitted model}
% =============================================================================

\begin{table}[H]
\centering
\begin{tabular}{@{}lrrrr@{}}
\toprule
\textbf{Parameter} & \textbf{Mean} & \textbf{SD} & \textbf{2.5\%} & \textbf{97.5\%} \\
\midrule
$\beta_0$ population baseline ratio      & 2.2920 & 0.0669 & 2.1577 & 2.4272 \\
$\beta_1$ population progression, /yr    & $-0.0927$ & 0.0121 & $-0.1159$ & $-0.0682$ \\
$\sigma$ total residual SD               & 1.1380 & 0.0203 & 1.0993 & 1.1778 \\
$\tau_0$ SD of patient baselines         & 1.3408 & 0.0499 & 1.2466 & 1.4421 \\
$\tau_1$ SD of patient slopes, /yr       & 0.1033 & 0.0137 & 0.0780 & 0.1315 \\
\bottomrule
\end{tabular}
\caption{Posterior summary, 2{,}000 retained draws. First and second halves of
the chain agree to four decimals on every parameter.}
\end{table}

\begin{figure}[H]
\centering
\includegraphics[width=\textwidth]{growth_trajectories.png}
\caption{The fitted model. Left: the population trajectory with its credible
interval, and 60 patients drawn from the random-effect distribution. Right:
the posterior for the population progression rate.}
\end{figure}

\noindent
\textbf{The population trajectory declines,} at 0.0927 per year with
$P(\beta_1 < 0) = 1.000$. Starting from a mean baseline ratio of 2.292, the
population average reaches the TD threshold after a median \textbf{12.3 years}
(95\% CI 9.9 to 16.4) and the PIGD threshold after 15.0 years (12.1 to 19.9).

\textbf{Patients differ substantially.} The between-patient SD of the slope is
0.103 per year, comparable in size to the population slope itself. A patient
one standard deviation faster than average progresses at roughly $-0.196$ per
year, reaching the TD threshold in about six years rather than twelve.

% =============================================================================
\section{Head-to-head with the hidden Markov model}
% =============================================================================

Both models are fitted on the same 755 training patients and scored on the same
held-out transitions --- those where every visit involved has a finite ratio.

\begin{table}[H]
\centering
\begin{tabular}{@{}lrrr@{}}
\toprule
\textbf{Model} & \textbf{Log-lik / transition} & \textbf{Accuracy} & \textbf{Brier} \\
\midrule
HMM (latent states)          & $\mathbf{-0.6559}$ & \textbf{0.761} & \textbf{0.3573} \\
Hierarchical Bayesian growth & $-0.6814$ & 0.737 & 0.3786 \\
\bottomrule
\end{tabular}
\caption{825 shared held-out transitions.}
\end{table}

\begin{table}[H]
\centering
\begin{tabular}{@{}rrrrr@{}}
\toprule
\textbf{Split} & \textbf{$n$} & \textbf{HMM} & \textbf{HB growth} & \textbf{Difference} \\
\midrule
0 & 761 & $-0.6948$ & $-0.7350$ & $+0.0403$ \\
1 & 832 & $-0.6822$ & $-0.7544$ & $+0.0722$ \\
2 & 817 & $-0.7111$ & $-0.7897$ & $+0.0785$ \\
3 & 778 & $-0.7264$ & $-0.7731$ & $+0.0467$ \\
4 & 748 & $-0.7028$ & $-0.7760$ & $+0.0732$ \\
5 & 786 & $-0.7355$ & $-0.7717$ & $+0.0362$ \\
\bottomrule
\end{tabular}
\caption{Six independent patient-level splits, both models refitted each time.
The HMM is ahead in 6 of 6, mean $+0.0579$ (sd 0.0188). Paired $t$-test
$p = 6.6\times10^{-4}$; Wilcoxon $p = 0.031$.}
\end{table}

\begin{figure}[H]
\centering
\includegraphics[width=0.86\textwidth]{hb_vs_hmm.png}
\caption{The same comparison. Higher is better.}
\end{figure}

\subsection{Why the growth model loses at this task}

\begin{keypoint}
The total residual SD is $\sigma = 1.138$. The Indeterminate band is
$1.15 - 0.90 = 0.25$ wide. \textbf{The noise is 4.6 times the width of the band
the model must predict membership of.}

Even with a perfectly estimated trajectory, the next observation lands almost
anywhere, so the Indeterminate class is nearly unhittable. The Markov model
does not face this: it predicts the class directly and never has to locate a
continuous value inside a narrow interval.
\end{keypoint}

\noindent
A second symptom of the same mismatch: a Normal observation model on the raw
ratio places \textbf{12.2\% of its five-year predictive mass below zero}, and a
ratio of two non-negative scores cannot be negative. The ratio is also strongly
right-skewed (skew 1.27, range 0 to 11.36), which a symmetric Normal cannot
represent.

\subsection{The two clinical tables side by side}

\begin{table}[H]
\centering
\begin{tabular}{@{}lrrr|rrr@{}}
\toprule
& \multicolumn{3}{c|}{\textbf{HMM (latent states)}} & \multicolumn{3}{c}{\textbf{HB growth}} \\
\textbf{Horizon} & TD & Ind & PIGD & TD & Ind & PIGD \\
\midrule
1 year  & 0.910 & 0.071 & 0.019 & 0.759 & 0.052 & 0.189 \\
2 years & 0.834 & 0.111 & 0.055 & 0.741 & 0.052 & 0.206 \\
3 years & 0.769 & 0.133 & 0.099 & 0.714 & 0.061 & 0.224 \\
5 years & 0.660 & 0.149 & 0.191 & 0.680 & 0.052 & 0.268 \\
\bottomrule
\end{tabular}
\caption{For a patient currently TD (HMM) or at the population baseline ratio of
2.29 (growth model).}
\end{table}

\noindent
The growth model gives a patient whose ratio is 2.29 --- twice the TD threshold
--- an 18.9\% chance of being PIGD within one year. That is not clinically
credible, and it is the $\sigma$ problem showing up directly in the deliverable.
Note also that its probabilities barely move with the horizon, because the
noise dominates the trend over five years.

% =============================================================================
\section{What the growth model provides that the Markov model cannot}
% =============================================================================

The comparison above is only about one task. On everything else the growth
model is the more useful object.

\begin{table}[H]
\centering
\begin{tabular}{@{}p{7.3cm}cc@{}}
\toprule
\textbf{Quantity} & \textbf{HMM} & \textbf{HB growth} \\
\midrule
Predicts the next recorded subtype well          & \textbf{yes} & weaker \\
Progression rate on the clinical scale, with CI  & no  & \textbf{yes} \\
Patient-specific progression rate                & no  & \textbf{yes} \\
Estimated years to cross a threshold             & no  & \textbf{yes} \\
Uses visits with an undefined ratio (16.3\%)     & \textbf{yes} & no \\
Quantifies label misclassification               & \textbf{yes} & no \\
Extends to the $\sigma_i$ stability phenotype    & no  & \textbf{yes} \\
Accepts baseline covariates naturally            & awkward & \textbf{yes} \\
\bottomrule
\end{tabular}
\end{table}

\begin{keypoint}
The Markov model reports \emph{that} a patient is likely to change state. The
growth model reports \emph{how fast their phenotype is moving}, on the clinical
scale, with uncertainty --- and lets that rate differ between patients. The
second half of the proposal, in which patient-specific variability $\sigma_i$
becomes a phenotype in its own right, is only expressible in the growth model.
\end{keypoint}

% =============================================================================
\section{Recommendation}
% =============================================================================

\begin{enumerate}
  \item \textbf{Keep both.} They are not competing answers to one question. The
        Markov model is the better classifier of the next visit; the growth
        model is the better description of disease progression.

  \item \textbf{Model the ratio on a transformed scale.} The three symptoms
        above --- $\sigma$ far exceeding the band width, negative predictions,
        strong right skew --- all follow from putting a symmetric Normal on a
        non-negative, skewed quantity. A log or similar transform would address
        all three at once, and the thresholds carry over exactly, since
        $R > 1.15 \iff \log R > \log 1.15$. This needs care because the ratio is
        exactly zero at some visits, so a shifted transform or a censored
        treatment is required.

  \item \textbf{Resolve the two-variance question deliberately.} Either adopt
        the state-space form, where the two noise sources are distinguishable,
        or state in the manuscript that only the total is identified. The
        current specification cannot support reporting them separately.

  \item \textbf{Model the two subscores rather than their ratio.} This would
        recover the 16.3\% of visits that have no finite ratio, and would remove
        the zero-denominator problem at its source.

  \item \textbf{Then add covariates and $\sigma_i$.} Once the observation scale
        is settled, the second half of the proposal follows directly.
\end{enumerate}

\vspace{0.8em}
\noindent\textbf{Reproducing this report.}
\begin{verbatim}
    python scripts/03_progression_model.py     # HMM and Markov baselines
    python scripts/04_hierarchical_bayes.py    # growth model and comparison
    python -m pytest tests/ -q                 # 134 tests
\end{verbatim}
\noindent
Stage 4 writes to \code{outputs/hbayes/}. Stage 3 was not modified.

\end{document}
